\documentclass[11pt]{article}
\usepackage[margin=1in]{geometry}
\usepackage{amsmath,amssymb}
\usepackage{algorithm}
\usepackage{algpseudocode}
\usepackage{booktabs}
\usepackage{hyperref}

\title{Parallel Cyclic Reduction (PCR) for Tridiagonal Systems}
\author{Notes for \texttt{proj\_tridiag}}
\date{}

\begin{document}
\maketitle

\section{The Tridiagonal System}

We wish to solve $Ax = d$ where $A$ is an $n \times n$ tridiagonal matrix:
\[
\begin{pmatrix}
b_0 & c_0 \\
a_1 & b_1 & c_1 \\
    & a_2 & b_2 & c_2 \\
    &     & \ddots & \ddots & \ddots \\
    &     &        & a_{n-2} & b_{n-2} & c_{n-2} \\
    &     &        &         & a_{n-1} & b_{n-1}
\end{pmatrix}
\begin{pmatrix}
x_0 \\ x_1 \\ x_2 \\ \vdots \\ x_{n-2} \\ x_{n-1}
\end{pmatrix}
=
\begin{pmatrix}
d_0 \\ d_1 \\ d_2 \\ \vdots \\ d_{n-2} \\ d_{n-1}
\end{pmatrix}
\]

Each row $i$ of this system gives the equation:
\begin{equation}\label{eq:row}
a_i \, x_{i-1} + b_i \, x_i + c_i \, x_{i+1} = d_i,
\end{equation}
with the convention $a_0 = 0$ and $c_{n-1} = 0$.

\section{The Key Idea: Eliminating Neighbours}

The Thomas algorithm (sequential LU) eliminates unknowns one at a time,
producing a serial dependency chain. PCR instead eliminates \emph{both}
neighbours $x_{i-1}$ and $x_{i+1}$ simultaneously from every equation,
at the cost of coupling to $x_{i-2}$ and $x_{i+2}$. After $\lceil \log_2 n \rceil$
rounds every equation depends only on its own unknown, yielding the solution directly.

\subsection{Derivation of one elimination step}

Consider three consecutive equations at step $s$. We use superscripts to
denote the step: $(a_i^{(s)}, b_i^{(s)}, c_i^{(s)}, d_i^{(s)})$.

At step $s$, every row $i$ satisfies:
\begin{equation}\label{eq:step}
a_i^{(s)} \, x_{i - 2^{s-1}} \;+\; b_i^{(s)} \, x_i \;+\; c_i^{(s)} \, x_{i + 2^{s-1}} \;=\; d_i^{(s)}.
\end{equation}

Initially ($s=0$), each equation couples $x_{i-1}, x_i, x_{i+1}$.
We want to eliminate $x_{i-2^{s-1}}$ and $x_{i+2^{s-1}}$ from equation $i$
by using the equations at rows $i - 2^{s-1}$ and $i + 2^{s-1}$.

\paragraph{Eliminating $x_{i-2^{s-1}}$.}
Row $i - 2^{s-1}$ at step $s$ reads:
\[
a_{i-2^{s-1}}^{(s)} \, x_{i - 2^s} + b_{i-2^{s-1}}^{(s)} \, x_{i-2^{s-1}} + c_{i-2^{s-1}}^{(s)} \, x_i = d_{i-2^{s-1}}^{(s)}.
\]
Solving for $x_{i-2^{s-1}}$:
\[
x_{i-2^{s-1}} = \frac{1}{b_{i-2^{s-1}}^{(s)}} \Bigl( d_{i-2^{s-1}}^{(s)} - a_{i-2^{s-1}}^{(s)} \, x_{i-2^s} - c_{i-2^{s-1}}^{(s)} \, x_i \Bigr).
\]
Define the multiplier:
\[
\alpha_i = -\frac{a_i^{(s)}}{b_{i-2^{s-1}}^{(s)}}.
\]
Multiplying row $i - 2^{s-1}$ by $\alpha_i$ and adding to row $i$ eliminates $x_{i-2^{s-1}}$.

\paragraph{Eliminating $x_{i+2^{s-1}}$.}
Similarly, from row $i + 2^{s-1}$:
\[
a_{i+2^{s-1}}^{(s)} \, x_i + b_{i+2^{s-1}}^{(s)} \, x_{i+2^{s-1}} + c_{i+2^{s-1}}^{(s)} \, x_{i+2^s} = d_{i+2^{s-1}}^{(s)}.
\]
Define:
\[
\gamma_i = -\frac{c_i^{(s)}}{b_{i+2^{s-1}}^{(s)}}.
\]
Multiplying row $i + 2^{s-1}$ by $\gamma_i$ and adding to row $i$ eliminates $x_{i+2^{s-1}}$.

\paragraph{Combined update.}
After both eliminations, the new coefficients for row $i$ at step $s+1$ are:
\begin{align}
a_i^{(s+1)} &= \alpha_i \cdot a_{i-2^{s-1}}^{(s)}, \label{eq:update_a} \\
b_i^{(s+1)} &= b_i^{(s)} + \alpha_i \cdot c_{i-2^{s-1}}^{(s)} + \gamma_i \cdot a_{i+2^{s-1}}^{(s)}, \label{eq:update_b} \\
c_i^{(s+1)} &= \gamma_i \cdot c_{i+2^{s-1}}^{(s)}, \label{eq:update_c} \\
d_i^{(s+1)} &= d_i^{(s)} + \alpha_i \cdot d_{i-2^{s-1}}^{(s)} + \gamma_i \cdot d_{i+2^{s-1}}^{(s)}. \label{eq:update_d}
\end{align}

After this step, row $i$ now couples unknowns at distance $2^s$ instead of $2^{s-1}$.

\paragraph{Boundary handling.}
When an index falls outside $[0, n)$, the corresponding neighbour does not exist.
We handle this by setting:
\[
a_i^{(s)} = 0 \;\text{if } i - 2^{s-1} < 0, \qquad
c_i^{(s)} = 0 \;\text{if } i + 2^{s-1} \geq n.
\]
This is equivalent to saying the multiplier $\alpha_i = 0$ or $\gamma_i = 0$
respectively, so no elimination is performed in that direction.

\section{Convergence}

At step $s$, each equation $i$ depends on $x_{i-2^{s-1}}$, $x_i$, and $x_{i+2^{s-1}}$.
After $S = \lceil \log_2 n \rceil$ steps, we have $2^{S-1} \geq n/2$, so
$i - 2^{S-1} < 0$ and $i + 2^{S-1} \geq n$ for every $i$. Therefore
$a_i^{(S)} = c_i^{(S)} = 0$ and each equation reduces to:
\[
b_i^{(S)} \, x_i = d_i^{(S)} \quad \Longrightarrow \quad x_i = \frac{d_i^{(S)}}{b_i^{(S)}}.
\]

\section{Explicit Example: $n = 4$}

Consider the system with $a_0 = 0$, $c_3 = 0$:
\[
\begin{pmatrix}
b_0 & c_0 &     &     \\
a_1 & b_1 & c_1 &     \\
    & a_2 & b_2 & c_2 \\
    &     & a_3 & b_3
\end{pmatrix}
\begin{pmatrix} x_0 \\ x_1 \\ x_2 \\ x_3 \end{pmatrix}
=
\begin{pmatrix} d_0 \\ d_1 \\ d_2 \\ d_3 \end{pmatrix}.
\]

\subsection{Step $s = 1$ (stride $2^0 = 1$)}

Each row $i$ currently couples $(x_{i-1}, x_i, x_{i+1})$.
We eliminate $x_{i-1}$ and $x_{i+1}$ using neighbours at distance 1.

\paragraph{Row 0:} $i - 1 = -1 < 0$, so $\alpha_0 = 0$. $i + 1 = 1$, so:
\[
\gamma_0 = -\frac{c_0}{b_1}, \qquad
a_0' = 0, \quad
b_0' = b_0 + \gamma_0 \, a_1, \quad
c_0' = \gamma_0 \, c_1, \quad
d_0' = d_0 + \gamma_0 \, d_1.
\]

\paragraph{Row 1:} Both neighbours exist:
\[
\alpha_1 = -\frac{a_1}{b_0}, \quad \gamma_1 = -\frac{c_1}{b_2}.
\]
\[
a_1' = \alpha_1 \, a_0 = 0, \quad
b_1' = b_1 + \alpha_1 \, c_0 + \gamma_1 \, a_2, \quad
c_1' = \gamma_1 \, c_2, \quad
d_1' = d_1 + \alpha_1 \, d_0 + \gamma_1 \, d_2.
\]
Note $a_1' = 0$ because $a_0 = 0$.

\paragraph{Row 2:} Both neighbours exist:
\[
\alpha_2 = -\frac{a_2}{b_1}, \quad \gamma_2 = -\frac{c_2}{b_3}.
\]
\[
a_2' = \alpha_2 \, a_1, \quad
b_2' = b_2 + \alpha_2 \, c_1 + \gamma_2 \, a_3, \quad
c_2' = \gamma_2 \, c_3 = 0, \quad
d_2' = d_2 + \alpha_2 \, d_1 + \gamma_2 \, d_3.
\]

\paragraph{Row 3:} $i + 1 = 4 \geq n$, so $\gamma_3 = 0$:
\[
\alpha_3 = -\frac{a_3}{b_2}, \qquad
a_3' = \alpha_3 \, a_2, \quad
b_3' = b_3 + \alpha_3 \, c_2, \quad
c_3' = 0, \quad
d_3' = d_3 + \alpha_3 \, d_2.
\]

After step 1, each row $i$ couples $(x_{i-2}, x_i, x_{i+2})$:
\[
\begin{aligned}
&\text{Row 0:} & &b_0' \, x_0 + c_0' \, x_2 = d_0', \\
&\text{Row 1:} & &b_1' \, x_1 + c_1' \, x_3 = d_1', \\
&\text{Row 2:} & a_2' \, x_0 + &b_2' \, x_2 = d_2', \\
&\text{Row 3:} & a_3' \, x_1 + &b_3' \, x_3 = d_3'.
\end{aligned}
\]

\subsection{Step $s = 2$ (stride $2^1 = 2$)}

Now we eliminate the remaining couplings at distance 2.

\paragraph{Row 0:} $i - 2 = -2 < 0$, and $i + 2 = 2$:
\[
\gamma_0'' = -\frac{c_0'}{b_2'}, \qquad
b_0'' = b_0' + \gamma_0'' \, a_2', \quad
d_0'' = d_0' + \gamma_0'' \, d_2'.
\]
Now $a_0'' = 0$, $c_0'' = 0$ (since $c_2' = 0$). So: $x_0 = d_0'' / b_0''$.

\paragraph{Row 1:} $i - 2 = -1 < 0$, and $i + 2 = 3$:
\[
\gamma_1'' = -\frac{c_1'}{b_3'}, \qquad
b_1'' = b_1' + \gamma_1'' \, a_3', \quad
d_1'' = d_1' + \gamma_1'' \, d_3'.
\]
So: $x_1 = d_1'' / b_1''$.

\paragraph{Row 2:} $i - 2 = 0$ exists, $i + 2 = 4 \geq n$:
\[
\alpha_2'' = -\frac{a_2'}{b_0'}, \qquad
b_2'' = b_2' + \alpha_2'' \, c_0', \quad
d_2'' = d_2' + \alpha_2'' \, d_0'.
\]
So: $x_2 = d_2'' / b_2''$.

\paragraph{Row 3:} $i - 2 = 1$ exists, $i + 2 = 5 \geq n$:
\[
\alpha_3'' = -\frac{a_3'}{b_1'}, \qquad
b_3'' = b_3' + \alpha_3'' \, c_1', \quad
d_3'' = d_3' + \alpha_3'' \, d_1'.
\]
So: $x_3 = d_3'' / b_3''$.

After $S = \lceil \log_2 4 \rceil = 2$ steps, every equation is decoupled, and we
read off each $x_i$ by simple division.

\section{Numerical Example}

Let us trace through with concrete numbers:
\[
A = \begin{pmatrix}
1 & 1 \\
2.1 & 2 & 1.1 \\
 & 4.1 & 3 & 1.2 \\
 & & 1.9 & 4
\end{pmatrix}, \quad
d = \begin{pmatrix} 0.1 \\ 0.5 \\ 0.3 \\ 0.2 \end{pmatrix}.
\]

So: $a = (0, 2.1, 4.1, 1.9)$, $b = (1, 2, 3, 4)$, $c = (1, 1.1, 1.2, 0)$.

\paragraph{Step 1.}
\begin{align*}
\text{Row 0:}\quad & \alpha_0 = 0, \quad \gamma_0 = -c_0/b_1 = -1/2 = -0.5 \\
& b_0' = 1 + (-0.5)(2.1) = -0.05, \quad c_0' = (-0.5)(1.1) = -0.55 \\
& d_0' = 0.1 + (-0.5)(0.5) = -0.15 \\[6pt]
\text{Row 1:}\quad & \alpha_1 = -a_1/b_0 = -2.1/1 = -2.1, \quad \gamma_1 = -c_1/b_2 = -1.1/3 \approx -0.3\overline{6} \\
& a_1' = (-2.1)(0) = 0, \quad c_1' = (-0.3\overline{6})(1.2) = -0.44 \\
& b_1' = 2 + (-2.1)(1) + (-0.3\overline{6})(4.1) = 2 - 2.1 - 1.503\overline{3} = -1.603\overline{3} \\
& d_1' = 0.5 + (-2.1)(0.1) + (-0.3\overline{6})(0.3) = 0.5 - 0.21 - 0.11 = 0.18 \\[6pt]
\text{Row 2:}\quad & \alpha_2 = -a_2/b_1 = -4.1/2 = -2.05, \quad \gamma_2 = -c_2/b_3 = -1.2/4 = -0.3 \\
& a_2' = (-2.05)(2.1) = -4.305, \quad c_2' = (-0.3)(0) = 0 \\
& b_2' = 3 + (-2.05)(1.1) + (-0.3)(1.9) = 3 - 2.255 - 0.57 = 0.175 \\
& d_2' = 0.3 + (-2.05)(0.5) + (-0.3)(0.2) = 0.3 - 1.025 - 0.06 = -0.785 \\[6pt]
\text{Row 3:}\quad & \alpha_3 = -a_3/b_2 = -1.9/3 \approx -0.6\overline{3}, \quad \gamma_3 = 0 \\
& a_3' = (-0.6\overline{3})(4.1) = -2.596\overline{6}, \quad c_3' = 0 \\
& b_3' = 4 + (-0.6\overline{3})(1.2) = 4 - 0.76 = 3.24 \\
& d_3' = 0.2 + (-0.6\overline{3})(0.3) = 0.2 - 0.19 = 0.01
\end{align*}

\paragraph{Step 2.}
Now stride is 2. Each row eliminates coupling at distance 2.
\begin{align*}
\text{Row 0:}\quad & \gamma_0'' = -c_0'/b_2' = -(-0.55)/0.175 \approx 3.142857 \\
& b_0'' = -0.05 + (3.142857)(-4.305) \approx -0.05 - 13.53 = -13.58 \\
& d_0'' = -0.15 + (3.142857)(-0.785) \approx -0.15 - 2.467 = -2.617 \\
& x_0 = d_0''/b_0'' \approx -2.617 / -13.58 \approx 0.19271 \\[6pt]
\text{Row 1:}\quad & \gamma_1'' = -c_1'/b_3' = -(-0.44)/3.24 \approx 0.135802 \\
& b_1'' = -1.603\overline{3} + (0.135802)(-2.596\overline{6}) \approx -1.603 - 0.353 = -1.956 \\
& d_1'' = 0.18 + (0.135802)(0.01) \approx 0.18 + 0.00136 = 0.18136 \\
& x_1 = d_1''/b_1'' \approx 0.18136 / -1.956 \approx -0.09272 \\[6pt]
\text{Row 2:}\quad & \alpha_2'' = -a_2'/b_0' = -(-4.305)/(-0.05) = -86.1 \\
& b_2'' = 0.175 + (-86.1)(-0.55) = 0.175 + 47.355 = 47.53 \\
& d_2'' = -0.785 + (-86.1)(-0.15) = -0.785 + 12.915 = 12.13 \\
& x_2 = d_2''/b_2'' \approx 12.13/47.53 \approx 0.25521 \\[6pt]
\text{Row 3:}\quad & \alpha_3'' = -a_3'/b_1' = -(-2.596\overline{6})/(-1.603\overline{3}) \approx -1.6195 \\
& b_3'' = 3.24 + (-1.6195)(-0.44) = 3.24 + 0.7126 = 3.9526 \\
& d_3'' = 0.01 + (-1.6195)(0.18) = 0.01 - 0.2915 = -0.2815 \\
& x_3 = d_3''/b_3'' \approx -0.2815/3.9526 \approx -0.07123
\end{align*}

These values match the CPU Thomas algorithm output (and can be verified with
NumPy: \texttt{np.linalg.solve(A, d)}).

\section{Parallelism and Complexity}

\begin{itemize}
\item \textbf{Work:} Each step updates all $n$ equations using
  Eqs.~\eqref{eq:update_a}--\eqref{eq:update_d}, requiring $O(n)$ work per step.
  With $\lceil \log_2 n \rceil$ steps, total work is $O(n \log n)$.
\item \textbf{Span (depth):} All $n$ updates within a step are independent
  and can execute in parallel. So the parallel depth is $O(\log n)$.
\item \textbf{Comparison to Thomas:} Thomas is $O(n)$ work but $O(n)$ span
  (fully sequential). PCR trades extra work for exponentially less depth,
  making it ideal for GPUs.
\end{itemize}

\section{Double Buffering}

The update equations read from step $s$ and write to step $s+1$.
Since all rows are updated simultaneously, we cannot overwrite in place.
Two buffers suffice: read from buffer $A$, write to buffer $B$, then swap.
This is the ``double buffering'' or ``ping-pong'' technique used in
the CUDA implementation.

A \texttt{\_\_syncthreads()} barrier is required between steps to ensure
all threads have finished writing before any thread reads the next step's
input.

\section{Block-Stride Loops}

When $n > \texttt{blockDim.x}$, a single thread must process multiple
elements. The inner loops iterate as:
\begin{verbatim}
    for (int i = threadIdx.x; i < n; i += blockDim.x) {
        // process element i
    }
\end{verbatim}
This ``block-stride'' pattern ensures coalesced memory access and full
utilization of available threads regardless of problem size.

\section{Application: Cubic Spline Fitting}

A natural application of tridiagonal solvers---and the original motivation
for the GPU PCR implementation---is cubic spline interpolation.

\subsection{Setup}

Given data points $(x_0, y_0),\, (x_1, y_1),\, \ldots,\, (x_n, y_n)$ with
$x_0 < x_1 < \cdots < x_n$, a cubic spline $S$ is a piecewise cubic
polynomial satisfying:
\begin{enumerate}
\item \textbf{Interpolation:} $S(x_i) = y_i$ for all $i$.
\item \textbf{$C^1$ continuity:} $S'$ is continuous at each interior knot.
\item \textbf{$C^2$ continuity:} $S''$ is continuous at each interior knot.
\end{enumerate}
Let $h_i = x_{i+1} - x_i$ denote the interval widths, and let
$z_i = S''(x_i)$ be the unknown second derivatives at the knots.

\subsection{Deriving the tridiagonal system}

On each interval $[x_i, x_{i+1}]$, $S$ is cubic, so $S''$ is linear.
Interpolating $S''$ between its endpoint values gives:
\[
S_i''(x) = z_i \,\frac{x_{i+1} - x}{h_i} + z_{i+1}\,\frac{x - x_i}{h_i}.
\]
Integrating twice and applying $S_i(x_i) = y_i$, $S_i(x_{i+1}) = y_{i+1}$
yields:
\begin{align*}
S_i(x) &= \frac{z_i}{6h_i}(x_{i+1}-x)^3
         + \frac{z_{i+1}}{6h_i}(x - x_i)^3 \\
       &\quad + \left(\frac{y_i}{h_i} - \frac{z_i\,h_i}{6}\right)(x_{i+1}-x)
         + \left(\frac{y_{i+1}}{h_i} - \frac{z_{i+1}\,h_i}{6}\right)(x - x_i).
\end{align*}
Differentiating, we obtain the first derivative at $x_{i+1}$ as
computed from the polynomial on its \emph{left} interval ($S_i$)
and its \emph{right} interval ($S_{i+1}$).
From the left:
\[
S_i'(x_{i+1}) = -\frac{z_i\,h_i}{6} + \frac{z_{i+1}\,h_i}{3}
  + \frac{y_{i+1} - y_i}{h_i},
\]
From the right:
\[
S_{i+1}'(x_{i+1}) = -\frac{z_{i+1}\,h_{i+1}}{3} + \frac{z_{i+2}\,h_{i+1}}{6}
  + \frac{y_{i+2} - y_{i+1}}{h_{i+1}}.
\]
The $C^1$ condition $S_i'(x_{i+1}) = S_{i+1}'(x_{i+1})$ at each
interior knot $i+1$ (for $i = 0, \ldots, n-2$) gives:
\begin{equation}\label{eq:spline_tridiag}
h_i \, z_i + 2(h_i + h_{i+1})\, z_{i+1} + h_{i+1}\, z_{i+2}
= 6\left(\frac{y_{i+2} - y_{i+1}}{h_{i+1}} - \frac{y_{i+1} - y_i}{h_i}\right).
\end{equation}
Re-indexing with $j = i+1$ for the interior knots $j = 1, \ldots, n-1$:
\[
\underbrace{h_{j-1}}_{a_j}\, z_{j-1}
+ \underbrace{2(h_{j-1} + h_j)}_{b_j}\, z_j
+ \underbrace{h_j}_{c_j}\, z_{j+1}
= \underbrace{6\!\left(\frac{y_{j+1} - y_j}{h_j}
  - \frac{y_j - y_{j-1}}{h_{j-1}}\right)}_{d_j}.
\]
This is a tridiagonal system of $n-1$ equations in $n+1$ unknowns
$(z_0, \ldots, z_n)$. Two boundary conditions close the system.

\subsection{Boundary conditions}

The choice of boundary condition determines the spline type
and how the first/last rows of the system are modified.

\begin{table}[h]
\centering
\renewcommand{\arraystretch}{1.8}
\begin{tabular}{l|c|c|c}
\toprule
 & $S(x)$ & $S'(x)$ & $S''(x)$ \\
\midrule
\textbf{Natural}
  & ---
  & ---
  & $S''(x_0) = 0,\; S''(x_n) = 0$ \\
\textbf{Clamped}
  & ---
  & $S'(x_0) = s_0,\; S'(x_n) = s_n$
  & --- \\
\textbf{Periodic}
  & ---\textsuperscript{*}
  & $S'(x_0) = S'(x_n)$
  & $S''(x_0) = S''(x_n)$ \\
\bottomrule
\end{tabular}
\caption{Boundary conditions by spline type. A ``---'' indicates no
extra constraint is imposed on that derivative order at the endpoints.\\
\textsuperscript{*}Periodic splines require $S(x_0) = S(x_n)$ (i.e.\ $y_0 = y_n$)
as a precondition on the input data, but this is not a boundary equation
in the tridiagonal system---it is already satisfied by the interpolation conditions.}
\label{tab:bc}
\end{table}

\paragraph{Natural spline.}
Set $z_0 = z_n = 0$ (zero curvature at endpoints). This reduces the system
to $n-1$ equations in $n-1$ unknowns $(z_1, \ldots, z_{n-1})$, a standard
tridiagonal system solved by Thomas or PCR.

\paragraph{Clamped spline.}
Prescribe the first derivatives $S'(x_0) = s_0$ and $S'(x_n) = s_n$.
The condition $S_0'(x_0) = s_0$ gives:
\[
\frac{h_0}{3}\,z_0 + \frac{h_0}{6}\,z_1
= \frac{y_1 - y_0}{h_0} - s_0,
\]
and $S_{n-1}'(x_n) = s_n$ gives:
\[
\frac{h_{n-1}}{6}\,z_{n-1} + \frac{h_{n-1}}{3}\,z_n
= s_n - \frac{y_n - y_{n-1}}{h_{n-1}}.
\]
These become the first and last rows, yielding an $(n+1)\times(n+1)$
tridiagonal system.

\paragraph{Periodic spline.}
Enforce $S'(x_0) = S'(x_n)$ and $S''(x_0) = S''(x_n)$, i.e.\ $z_0 = z_n$.
This wraps the tridiagonal system into a \textbf{cyclic tridiagonal} system,
which can be solved via the Sherman--Morrison formula (reducing it to two
standard tridiagonal solves).

\subsection{Sherman--Morrison for Cyclic Tridiagonals}

The periodic spline produces a \textbf{cyclic tridiagonal} system:
the matrix has two extra corner entries $\alpha$ (bottom-left) and
$\beta$ (top-right) in addition to the usual three diagonals:
\[
\tilde{A} =
\begin{pmatrix}
b_0    & c_0    &        &          & \beta  \\
a_1    & b_1    & c_1    &          &        \\
       & \ddots & \ddots & \ddots   &        \\
       &        & a_{m-2}& b_{m-2}  & c_{m-2}\\
\alpha &        &        & a_{m-1}  & b_{m-1}
\end{pmatrix}.
\]
We want to solve $\tilde{A}\,x = r$.
A standard tridiagonal solver cannot handle the corners directly.
The Sherman--Morrison formula lets us decompose this into two
standard tridiagonal solves plus a cheap combine step.

\paragraph{The Sherman--Morrison formula.}
Recall: given an invertible matrix $A'$ and vectors $u, v$ such that
$\tilde{A} = A' + u\,v^T$ (a \textbf{rank-1 update}), the solution
of $\tilde{A}\,x = r$ is:
\[
x = \underbrace{A'^{-1}\,r}_{y}
  \;-\; \frac{v^T\!\overbrace{A'^{-1}\,r}^{y}}
             {1 + v^T\!\underbrace{A'^{-1}\,u}_{z}}
       \;\cdot\;\underbrace{A'^{-1}\,u}_{z}.
\]
Crucially, you never form $A'^{-1}$---you just solve two linear
systems with the \emph{same} coefficient matrix $A'$:
\begin{enumerate}
\item Solve $A'\,y = r$ \quad (standard tridiagonal solve),
\item Solve $A'\,z = u$ \quad (standard tridiagonal solve),
\end{enumerate}
then combine:
\[
x = y - \frac{v^T y}{1 + v^T z}\; z.
\]

\paragraph{Constructing the decomposition.}
Choose $\gamma = -b_0$ (any nonzero value works; this particular
choice ensures $b_0' = b_0 - \gamma = 2\,b_0$, avoiding catastrophic
cancellation that would occur if $\gamma \approx b_0$). It is also
\textbf{scale-invariant}: the modification is always a factor of~2
on $b_0$, regardless of its magnitude. A na\"ive choice like
$\gamma = -1$ would give $b_0' = b_0 + 1$, which distorts the
relative scale when $b_0 \ll 1$ or $b_0 \gg 1$, potentially
worsening the condition number of $A'$.

With this choice, define:
\[
u = \begin{pmatrix} \gamma \\ 0 \\ \vdots \\ 0 \\ \alpha \end{pmatrix},
\qquad
v = \begin{pmatrix} 1 \\ 0 \\ \vdots \\ 0 \\ \beta/\gamma \end{pmatrix}.
\]
The outer product $u\,v^T$ is a matrix with only four nonzero entries:
\[
(u\,v^T)_{ij} =
\begin{cases}
\gamma            & (i,j) = (0,\,0), \\
\beta             & (i,j) = (0,\,m{-}1), \\
\alpha            & (i,j) = (m{-}1,\,0), \\
\alpha\beta/\gamma & (i,j) = (m{-}1,\,m{-}1).
\end{cases}
\]
Setting $A' = \tilde{A} - u\,v^T$:
\begin{itemize}
\item The corners are cancelled:
  $A'_{0,\,m-1} = \beta - \beta = 0$ and
  $A'_{m-1,\,0} = \alpha - \alpha = 0$.
\item The diagonal is modified:
  $b_0' = b_0 - \gamma = 2\,b_0$ and
  $b_{m-1}' = b_{m-1} - \alpha\beta/\gamma$.
\item All other entries are unchanged.
\end{itemize}
Thus $A'$ is a \textbf{standard tridiagonal} matrix, and
$\tilde{A} = A' + u\,v^T$ exactly.

\paragraph{The combine step.}
After obtaining $y$ and $z$, the dot products are cheap since $v$
is sparse:
\[
v^T y = y_0 + \frac{\beta}{\gamma}\,y_{m-1},
\qquad
v^T z = z_0 + \frac{\beta}{\gamma}\,z_{m-1}.
\]
The final solution is:
\[
x_i = y_i - \frac{v^T y}{1 + v^T z}\; z_i,
\qquad i = 0, \ldots, m{-}1.
\]
This last step is trivially parallel (element-wise).

\paragraph{GPU strategy.}
The two tridiagonal solves share the same coefficient matrix $A'$
but have different right-hand sides ($r$ and $u$). They can be:
\begin{itemize}
\item batched as two rows in a single PCR kernel call, or
\item executed as two sequential PCR calls (same $A'$, different RHS).
\end{itemize}
The combine step is a simple element-wise kernel. In a batched
setting where many periodic splines are solved simultaneously,
each spline simply contributes two rows to the batch.

\subsection{Why PCR fits well}

\begin{itemize}
\item The spline tridiagonal system is \textbf{diagonally dominant}:
  $b_j = 2(h_{j-1} + h_j) > h_{j-1} + h_j \geq |a_j| + |c_j|$,
  ensuring numerical stability without pivoting.
\item For \textbf{batched} spline fitting (thousands of independent splines),
  each spline produces one small tridiagonal system. These map directly
  to the rows of the batched PCR kernel.
\item Typical spline sizes (tens to hundreds of knots) fit comfortably
  in shared memory, enabling the shmem PCR variant.
\end{itemize}

\section{Summary}

\begin{enumerate}
\item PCR eliminates both neighbours simultaneously from every equation.
\item After each step, coupling distance doubles: $1 \to 2 \to 4 \to \cdots$
\item After $\lceil \log_2 n \rceil$ steps, every equation is decoupled.
\item All updates within a step are independent $\Rightarrow$ GPU-parallel.
\item Trade-off: $O(n \log n)$ work vs.\ $O(\log n)$ depth.
\end{enumerate}

\end{document}
